% =============================================================================
% ANEXO I — Diagramas de Estado (máquinas de estado del dominio)
% =============================================================================
\chapter{Anexo I — Diagramas de Estado}
\label{cap:anexo_estados}

Este anexo modela el ciclo de vida de las entidades con estados explícitos definidos en el
esquema real (\textit{check constraints} y enumeraciones), conforme a las migraciones del
repositorio.\par

\section{Ciclo de Vida del Pipeline de Reclutamiento}
\label{sec:anexo_estado_pipeline}
La columna \comando{stage} de \comando{core.profile\_likes} admite los valores
\comando{saved}, \comando{review}, \comando{interview}, \comando{offer} y \comando{hired}
(\textit{check constraint}, V6). La Figura~\ref{fig:est_pipeline} modela las transiciones.

\begin{figure}[!ht]
    \centering
    \begin{tikzpicture}[font=\sffamily\scriptsize, >=Latex, node distance=1.9cm,
        st/.style={draw=colorPrimario, fill=headerTabla, rounded corners=6pt,
            minimum width=1.7cm, minimum height=0.7cm, align=center},
        ini/.style={circle, fill=colorPrimario, minimum size=3mm, inner sep=0},
        tr/.style={-{Latex[length=1.6mm]}, colorPrimario}]
        \node[ini] (i) {};
        \node[st, right=1.0cm of i] (s) {saved};
        \node[st, right=of s] (r) {review};
        \node[st, right=of r] (e) {interview};
        \node[st, right=of e] (o) {offer};
        \node[st, right=of o] (h) {hired};
        \draw[tr] (i) -- (s);
        \draw[tr] (s) -- (r); \draw[tr] (r) -- (e); \draw[tr] (e) -- (o); \draw[tr] (o) -- (h);
        \draw[tr] (r.south) to[bend left=30] node[below,font=\tiny]{descarta} (s.south);
        \draw[tr] (e.north) to[bend right=30] node[above,font=\tiny]{retrocede} (r.north);
    \end{tikzpicture}
    \caption{Máquina de estados del \comando{stage} de un candidato en el pipeline.}
    \label{fig:est_pipeline}
\end{figure}

\section{Salud de una Conexión Externa}
\label{sec:anexo_estado_conexion}
La columna \comando{api\_health} de \comando{core.profile\_connections} toma los valores
\comando{healthy}, \comando{degraded} y \comando{down}, y \comando{status} alterna entre
\comando{connected}, \comando{disconnected} y \comando{pending}. La Figura~\ref{fig:est_conn}
modela la salud de la conexión.

\begin{figure}[!ht]
    \centering
    \begin{tikzpicture}[font=\sffamily\scriptsize, >=Latex, node distance=2.4cm,
        st/.style={draw=colorPrimario, fill=headerTabla, rounded corners=6pt,
            minimum width=1.8cm, minimum height=0.7cm, align=center},
        tr/.style={-{Latex[length=1.6mm]}, colorPrimario}]
        \node[st] (hh) {healthy};
        \node[st, right=of hh] (dg) {degraded};
        \node[st, right=of dg] (dn) {down};
        \draw[tr] (hh.north) to[bend left=20] node[above,font=\tiny]{latencia/errores} (dg.north);
        \draw[tr] (dg.south) to[bend left=20] node[below,font=\tiny]{recupera} (hh.south);
        \draw[tr] (dg.north) to[bend left=20] node[above,font=\tiny]{falla sync} (dn.north);
        \draw[tr] (dn.south) to[bend left=20] node[below,font=\tiny]{reconecta} (dg.south);
    \end{tikzpicture}
    \caption{Máquina de estados de \comando{api\_health} de una conexión externa.}
    \label{fig:est_conn}
\end{figure}

\section{Ciclo de Vida de un Código OTP}
\label{sec:anexo_estado_otp}
Los registros de \comando{core.profile\_security\_codes} nacen activos, expiran a los 15 minutos
o se consumen (\comando{used = true}). La Figura~\ref{fig:est_otp} lo modela.

\begin{figure}[!ht]
    \centering
    \begin{tikzpicture}[font=\sffamily\scriptsize, >=Latex, node distance=2.6cm,
        st/.style={draw=colorPrimario, fill=headerTabla, rounded corners=6pt,
            minimum width=1.9cm, minimum height=0.7cm, align=center},
        fin/.style={draw=grisISO, fill=grisTecnico, rounded corners=6pt,
            minimum width=1.9cm, minimum height=0.7cm, align=center},
        ini/.style={circle, fill=colorPrimario, minimum size=3mm, inner sep=0},
        tr/.style={-{Latex[length=1.6mm]}, colorPrimario}]
        \node[ini] (i) {};
        \node[st, right=1.0cm of i] (a) {emitido};
        \node[fin, above right=0.6cm and 2.2cm of a] (u) {consumido};
        \node[fin, below right=0.6cm and 2.2cm of a] (x) {expirado};
        \draw[tr] (i) -- (a);
        \draw[tr] (a) -- node[above,font=\tiny,sloped]{verifica OK} (u);
        \draw[tr] (a) -- node[below,font=\tiny,sloped]{15 min} (x);
    \end{tikzpicture}
    \caption{Ciclo de vida de un código OTP de seguridad.}
    \label{fig:est_otp}
\end{figure}

\section{Ciclo de Vida de la Cuenta y el Portafolio}
\label{sec:anexo_estado_cuenta}
Una cuenta transita de \comando{activa} a \comando{baja lógica} (\comando{deleted\_at}); el
portafolio alterna entre \comando{borrador}, \comando{publicado} y \comando{protegido} según
\comando{is\_published} y la contraseña. La Figura~\ref{fig:est_portfolio} modela el portafolio.

\begin{figure}[!ht]
    \centering
    \begin{tikzpicture}[font=\sffamily\scriptsize, >=Latex, node distance=2.6cm,
        st/.style={draw=colorPrimario, fill=headerTabla, rounded corners=6pt,
            minimum width=1.9cm, minimum height=0.7cm, align=center},
        tr/.style={-{Latex[length=1.6mm]}, colorPrimario}]
        \node[st] (b) {borrador};
        \node[st, right=of b] (p) {publicado};
        \node[st, right=of p] (pr) {protegido};
        \draw[tr] (b) -- node[above,font=\tiny]{publica} (p);
        \draw[tr] (p.south) to[bend left=25] node[below,font=\tiny]{despublica} (b.south);
        \draw[tr] (p) -- node[above,font=\tiny]{+contraseña} (pr);
        \draw[tr] (pr.south) to[bend left=25] node[below,font=\tiny]{quita clave} (p.south);
    \end{tikzpicture}
    \caption{Máquina de estados de la visibilidad del portafolio.}
    \label{fig:est_portfolio}
\end{figure}

\clearpage
